% !TEX root = chapter-standalone.tex

\section{Semifunctors}

\linkvideo{spring2021-functors:semi-and-fun:semi-fun-def} % Definition of semi-functor

%\linkvideo{spring2021-functors:semi-and-fun} % Semifunctors and functors

What is an appropriate analogue of `functor' for the world of semicategories?
One option is to simply drop the last condition in the definition of functor. We call the gadget we obtain in this way a \emph{semifunctor}. It is a generalization of the notion of semigroup homomorphism. 

\begin{ctdefinition}[Semifunctor]
    \label{def:semifunctor}
    Given \SY{semicategories}~\CatC and~\CatD, a \maindef{semifunctor}~$\funa\colon \CatC\fto \CatD$ from~\CatC to~\CatD is:

    \constit

    \begin{enumerate}
        \item A function
              \begin{equation}
                  \funaob \colon \ObC \to \ObD.
              \end{equation}
        \item For every pair of objects~$\Obja, \Objb\setin \Obof\CatC$, a function
              \begin{equation}
                  \funamor\colon \HomSet{\CatC}{\Obja}{\Objb} \to \HomSet{\CatD}{\funaob(\Obja)}{\funaob(\Objb)}.
              \end{equation}
    \end{enumerate}

    \condit

    \begin{enumerate}
        \item The function $\funamor$ is compatible with the composition operations in the source and target category, respectively:
              \begin{equation}
                  \prfperiod{
                      \mora\colon \Obja \mtoin{\CatC} \Objb
                  }{\quad}{
                      \morb\colon \Objb \mtoin{\CatC} \Objc
                  }{
                      \funamor(\mora \mthenof{\CatC} \morb)=\funamor(\mora) \mthenof{\CatD} \funamor(\morb)
                  }
              \end{equation}
    \end{enumerate}
\end{ctdefinition}

This situation is depicted graphically in \cref{fig:functor_detail}.
%
It is common to overload the notation and use~$\funa$ to denote not only the whole functor, both also both of it's constituent functions~$\funaob$ and~$\funamor$.
The diagram with this overloaded ``synthetic notation'' is in \cref{fig:functor_representation}.

\vfill

\begin{figure*}[h!]
    % \begin{ctdefinitionshade}
    \subfloat[\label{fig:functor_detail}
        Functor diagram] {
        \centering
        \includesag{095_functor_detail}
    }
    \subfloat[\label{fig:functor_representation}Synthetic notation] {
        \centering
        \includesag{095_functor}
    }
    % \end{ctdefinitionshade}
    \caption{
        Commuting diagrams for \SY{semifunctors}, with verbose notation (left) and synthetic notation (right).
    }
\end{figure*}

\vspace{2cm}
